\clearpage\thispagestyle{empty}
\pdfbookmark[0]{Quellen zu Gauß an Bode (9. August 1815)}{bode1815aug-sources}
\begin{center}\textbf{Quellen zu Gauß an Bode (9. August 1815)}\end{center}
\noindent Diese KI-Übertragung ins heutige Deutsch folgt dem Brief von Carl Friedrich Gauß an Johann Elert Bode in den \href{https://gdz.sub.uni-goettingen.de/id/PPN236007467}{\textit{Werken}, Band~VI} (Göttingen, 1874), S.~387--389. Bode druckte ihn unter einer Überschrift zu den Kometenbeobachtungen, den Bahnelementen und den Beobachtungen der Juno im \textit{Astronomischen Jahrbuch für 1818} (Berlin, 1815), S.~\href{https://api.digitale-sammlungen.de/iiif/image/v2/bsb10538347_00237/full/1600,/0/default.jpg}{229}--\href{https://api.digitale-sammlungen.de/iiif/image/v2/bsb10538347_00240/full/1600,/0/default.jpg}{232}. Die Links führen zum Exemplar der Bayerischen Staatsbibliothek.
\par\smallskip\noindent Das Jahrbuch gibt für die Rektaszension des Kometen am 30.~März 59~13~3 an (S.~\href{https://api.digitale-sammlungen.de/iiif/image/v2/bsb10538347_00238/full/1600,/0/default.jpg}{230}), die Werkausgabe im Brief 59~13~2. Auf S.~\href{https://api.digitale-sammlungen.de/iiif/image/v2/bsb10538347_00239/full/1600,/0/default.jpg}{231} steht im Jahrbuch \textit{Elemente}, in der Werkausgabe \textit{Elemnnte}. Die Deklination am 27.~Juli endet in beiden Drucken nach 20~37 mit einem Doppelpunkt.
\par\smallskip\noindent Auf S.~\href{https://api.digitale-sammlungen.de/iiif/image/v2/bsb10538347_00240/full/1600,/0/default.jpg}{232} nennt das Jahrbuch $77\tfrac13$ Jahre Umlaufzeit und für die Deklination der Juno am 29.~März $1\,32\,30{,}1$ N. In der Werkausgabe stehen $77\tfrac12$ Jahre und $1\,32\,38.1$ N. Bei den elliptischen Elementen steht im Jahrbuch „Abstand in der Sonnennähe“, und der aufsteigende Knoten hat keine Einheitenzeichen. Die Werkausgabe nennt den „Abstand von der Sonnennähe“ und gibt für den Knoten $83^\circ\,26'\,21''$ an. Beide Drucke nennen für die Montierung eine halbe Bogenminute. Die Schreibung \textit{Frauenhofer} folgt der Werkausgabe. KI-Übertragung und KI-Satz.
\clearpage
